
%
% Copyright Ian G. Abel 2018
% 

\documentclass[prb,aps,amssymb,amsmath,a4paper]{revtex4}
\usepackage{graphicx}
\usepackage{bm,ulem,color}

\sloppy

\include{gyrokinetics-macros}
\providecommand{\coll}[4]{ \ensuremath{C_{#1\,#2}[#3,#4]}}
\providecommand{\lorentz}[1]{\ensuremath{\mathcal{L}}[#1]}
\providecommand{\nuD}[2]{\ensuremath{\nu_{D}^{#1\,#2}}}
\providecommand{\nup}[2]{\ensuremath{\nu_{\parallel}^{#1\,#2}}}
\providecommand{\nus}[2]{\ensuremath{\nu_{s}^{#1\,#2}}}
\providecommand{\nu}[2]{\ensuremath{\nu_{#1\,#2}}}
\providecommand{\injeng}{\varepsilon_*}
\providecommand{\injvel}{v_*}
\providecommand{\grad}{\bm{\nabla}}
\begin{document}

\date{\today}

\title{The MEQ Algorithm}

\maketitle

\section{Conventional Formulation}
\begin{eqnarray}
\bm{q} &= \grad \psi / R \\
-\dv \bm{q}  &= F(R,z,\psi) / R\\
\psi &= \phi_h 
\end{eqnarray}
But we need to turn $\phi_h$ into a linear operator on $\bm{q}$ for implementation purposes.

Really $\phi_h$ is a vector of values at the nodes on the edge element. However, at a given node $\bm{x}$
\begin{equation}
\phi_h(\bm{x}) = - \int_0^{l(\bm{x})} R(s) \bm{q}( \bm{x} + \bm{t}(x) s )\cdot\bm{t}(\bm{x}) ds.
\end{equation}

So, for all points on an edge $\partial K$, we store an $l$ and a $\bm{t}$.
Using gaussian integration along the line,
		\begin{equation}
		\phi_h = - \sum_j w_j R(s_j/l) \bm{q}(\bm{x}_j)\cdot\bm{t} l
		\label{<++>}
		\end{equation}
		where now $s_j$ is on $(0,1)$. 

\end{document}

